\subsection{Further four- and five-column upper bounds}
\label{sec:further-finite}

\begin{theorem}\label{thm:further-finite}
The three parameters have the following common values:
\[
\begin{array}{c|ccccc}
(m,n)&(7,5)&(8,5)&(10,5)&(9,4)&(10,4)\\
\hline
\zr(m,n)=\zs(m,n)=\zt(m,n)&24&27&34&24&27.
\end{array}
\]
\end{theorem}

\begin{proof}
Since $\zr\le\zs\le\zt$, it suffices to give recursive-line lower
displays and upper bounds for $\zt$. The actual upper
bounds use the complete-matching method of this section: each excluded
local pattern has a shorter real SOS, while inherited bounds are applied
only to complete selected squares on a maximum base of the smaller
rectangle. No rejected recursive or signed certificate is an upper-bound
premise. Each target exceeds the asserted value by one:
\begin{center}
\begin{tabular}{ccrrrr}
\toprule
Rectangle&Maximum bases&$|F|$&Target pairs&Target holes&Target length\\
\midrule
$7\times5$&4&15&10&0&25\\
$8\times5$&1&17&11&1&28\\
$10\times5$&1&20&15&0&35\\
$9\times4$&3&15&10&1&25\\
$10\times4$&5&16&12&0&28\\
\bottomrule
\end{tabular}
\end{center}

Here ``maximum bases'' counts row/column isomorphism classes. Their
coverage is elementary. On five columns, pair counting gives the upper
bounds $15,17,20$. At seven rows, either two triple rows use all three
pairs in each triple, leaving four pair rows and one singleton row, or
one triple is accompanied by six of the seven remaining column pairs.
The triples in the first case meet in one column; the singleton is
anchored either there or in an outer column. In the second case the
omitted pair lies either outside the triple or crosses it. These are
four types, with 145 row multisets on labeled columns in total. At eight
rows the base is one triple and all seven remaining pair rows; at ten
rows it is all ten pair rows. Larger row degrees are excluded by the
same pair budget and $C_4$-freeness.

For four columns every maximum base consists of the six pair rows and
singleton rows. At nine rows the anchor multiplicities are
$(3,0,0,0)$, $(2,1,0,0)$ and $(1,1,1,0)$. At ten rows they are
$(4,0,0,0)$, $(3,1,0,0)$, $(2,2,0,0)$, $(2,1,1,0)$ and $(1,1,1,1)$. Thus the table covers every maximum base.

For each type, every pairing with the stated hole count either appears
in the matching domain or is excluded by a verified local reduction. The local exclusions are checked by
exact coefficient identities or certified lower-rank real Gram matrices.
The identities use integers, positive rational weights or $\Q(\sqrt3)$.
For the real Gram matrices, exact rational bounds keep the principal
block positive definite and make an invertible-preconditioned Schur map
a strict contraction of a closed box into itself. Its fixed point has
zero Schur complement, giving a positive-semidefinite matrix of the
smaller principal rank. The checks use exact LDL decompositions and
rational arithmetic; rationality of the final matrix entries is not assumed. The maximum-retaining seven-row, five-column deletion clauses use
$\zt(6,5)=20$ (Theorem~\ref{thm:small}, transposed).
The four-column inputs also use $\zt(6,4)=16$ from
\cite[Corollary~4.23]{LQ} and $\zt(7,4)=19$, $\zt(8,4)=21$ from
\cite[Theorems~4.2 and~4.4]{XY}; the ten-row input uses the nine-row
bound only after that separate proof is complete. Only the $10\times5$ formula uses lexicographic constraints: all $120$
column permutations, with their induced permutations of the pair rows,
act on its matching variables. A hypothetical irreducible display and
every image under this group satisfy every matching clause and every
reduction clause: each reduction is certified independently of the lex
constraints. The lexicographically least image therefore satisfies both the
reduction clauses and the inclusive lex comparisons; the reduction-clause
list need not itself be invariant.
Independent input reconstruction and
complete clausal refutations exclude every target in the table. The
lower certificates and corresponding exact upper-bound packets are
indexed in Appendix~\ref{app:data}. The four-column recursive values
were already known from~\cite{CC}; here the new conclusion is their
actual-SOS optimality.
\end{proof}

The value $\zt(10,5)=34$ also answers the five-column case of
\cite[Section~8, Open Problem~1]{CLQIncidence}.

